\subsection{Vector Floating-Point Element Types and Semantics}

\texttt{VECTORFP} requires \texttt{VECTOR} and \texttt{FP}. Its mnemonics
use the \texttt{VF} prefix and their element suffixes \texttt{H}, \texttt{S},
and \texttt{D} select IEEE 754 binary16, binary32, and binary64 elements.
Binary16 has default quiet NaN \texttt{0x7e00}. The three-bit \texttt{tzz}
selector values five through seven select \texttt{H/S/D}; value four remains
invalid. In FP-only families, the two-bit \texttt{zz} selector reserves zero
and assigns one through three to \texttt{H/S/D}.

\paragraph{Floating-Point Exceptions, Conversions, and Reductions.}

Each active FP lane applies the corresponding scalar rounding, NaN,
signed-zero, default-NaN, denormal, and exception rules.  Inactive lanes
generate no cause.  The instruction unions all active-lane causes before
testing the enables in the original \texttt{FSTATUS}.  If any generated cause
is enabled, one (:ref:FP.events.EXCEPTION.FLOATING_POINT_EXCEPTION:) occurs and neither a destination
effect nor newly accrued \texttt{FFLAGS} becomes visible.  Otherwise the union
is accrued once and all destination effects commit together.  FP-to-integer
conversion rounds toward zero and follows the scalar saturation and invalid
result rules.

FP reductions return an \texttt{Fn}.  Empty FP ADD returns positive zero, empty FP
MIN returns positive infinity, and empty FP MAX returns negative infinity.
FP ADD visits active lanes in increasing logical-lane order and may not be
reassociated.
